% Candidate rules: candidate guarantees (the four guarantees this project was asked to check).  Written by hand; see docs/one-id-survey.md.

\section{Candidate guarantees}

\begin{theorem}[Only requested and consented data is received (G-01)]
  \label{thm:g01}
  \lean{Rules.OneId.accepted_holds_only_requested}
  \uses{def:acceptance, thm:wallet-discloses, thm:no-check-shows-consent, thm:partial-consent, def:c26, def:c07, def:c12, def:c22, def:c58, def:c65, def:c75, def:element, thm:names-unique, def:c31}
  A relying party receives only the data elements it requested and the holder
  consented to.  Not proven.

  A wallet that follows the rules discloses only what was requested and
  approved, which is proven.  Whether an accepted response holds only what was
  requested depends on the relying party refusing extra data, which no text
  that was read requires.  And no check a relying party makes can show that
  the holder approved what was sent, so that half rests on the wallet.

  Counterexample: a check-in page requests only the family name.  A wallet
  that does not follow the rules also sends the portrait, with valid
  signatures made for this request.  A relying party that does not refuse
  extra data accepts the response and holds a face image it never asked for.

  Open point: OpenID4VP forbids a wallet to send a claim the query did not
  select and tells a verifier not to rely on the wallet for it, without
  saying that a response with more is refused (OpenID4VP §6.4, §14.9).  The
  mdoc profile lists no step that compares the response with the request
  (C-65).  Must a relying party refuse a response that holds a data element
  it did not request, or may it accept the response and discard the extra?
  And who vouches that a wallet asks for consent (C-01, C-08)?

  Source:
  \href{https://standards.atlassian.net/wiki/spaces/OneID/pages/2947416082/Technical+Implementation+Guide}{One ID guide} §4.1.1, §3.1, §4.2.1.3.
\end{theorem}

\begin{theorem}[An accepted presentation is bound to the device and the request (G-02)]
  \label{thm:g02}
  \lean{Rules.OneId.accepted_answers_this_request}
  \uses{def:acceptance, thm:accepted-device-key, thm:openid4vp-refuses-other-request, def:c78, def:c65, def:c72, def:c59, def:c88, def:proximity-session, def:c03, def:c24, def:c02}
  An accepted presentation is bound to the holder's device key and to the
  verifier's request: its nonce in the remote profiles, its session in
  proximity.  Not proven.

  The device half is proven for all three profiles: every accepted credential
  is signed with the device key its issuer signed.  The request half is proven
  for the OpenID4VP profile, where a response made for one nonce, origin and
  key is refused for any other.  It is open for the mdoc profile and for
  proximity.

  Counterexample: a response made for a request with one nonce is captured and
  sent again in answer to a later request with another nonce.  The OpenID4VP
  profile refuses it.  The mdoc profile accepts it if the relying party does
  not compare the session transcript with its own request.

  Open point: the OpenID4VP profile requires the relying party to rebuild the
  session transcript and validate the nonce.  The mdoc profile's list of
  checks has no such step (C-65), and the proximity profile lists the
  rebuilding with no requirement word (C-88).  Must a relying party in the
  mdoc profile and a reader in proximity check that a response was made for
  their own request?  The guarantee also rests on issuance for what a device
  key is bound to (C-03), and on the entropy of the nonce (C-72).

  Source:
  \href{https://standards.atlassian.net/wiki/spaces/OneID/pages/2947416082/Technical+Implementation+Guide}{One ID guide} §4.2.
\end{theorem}

\begin{theorem}[A photoid and a boarding pass presented together are one traveler's (G-03)]
  \label{thm:g03}
  \lean{Rules.OneId.accepted_pair_same_traveler}
  \uses{def:pair, thm:no-presentation-check, thm:matched-pair-same-traveler, thm:same-traveler-matches, def:c19, def:c42, def:c93, def:c94, def:c01, def:c02, def:c03, def:c104}
  A photoid and a boarding pass presented together refer to the same
  traveler, under a defined matching rule.  Not proven.

  It fails under every reading of the check: with nothing compared, with
  names compared, and with device keys compared.  Comparing device keys gives
  one traveler where issuance gives each device key to one traveler, which is
  proven under that condition and which the guide does not state.

  Counterexample: Alice Martin's photoid and Bruno Keller's boarding pass were
  issued to one shared phone.  Under the reading that the two credentials must
  carry the same device key, the pair is accepted: Alice's face goes with
  Bruno's boarding pass.

  Open point: a presentation gives a relying party one more thing to compare
  than the content of the two credentials: the device key each was issued to.
  That settles the question only if issuance makes a device key stand for one
  traveler, and the guide puts issuance out of scope (C-03).  What ties a
  photoid and a boarding pass presented together to one traveler: the same
  device key, a rule for issuers, a comparison of names, or something else,
  and what must a relying party check?  A boarding pass for an adult with an
  infant is a further case (C-104).

  Source:
  \href{https://standards.atlassian.net/wiki/spaces/OneID/pages/2947416082/Technical+Implementation+Guide}{One ID guide} §4.1.2, §3.2.
\end{theorem}

\begin{theorem}[An expired credential is never accepted (G-04)]
  \label{thm:g04}
  \lean{Rules.OneId.accepted_not_expired}
  \uses{def:acceptance, def:expired, def:c78, def:c101, thm:valid-at-validuntil, def:c02}
  A relying party never accepts a credential whose validity has ended.  Not
  proven.

  It is proven for the mdoc profile, where the guide says the clause it cites
  covers time validity.  It is open for the OpenID4VP profile and for
  proximity.  Within the guide's rules a credential from a trusted issuer
  ends only by expiry, since revocation is out of scope (C-02).

  Counterexample: a boarding pass whose validity ended at 10:00 is presented
  at 11:00 with valid signatures.  The mdoc profile refuses it.  The OpenID4VP
  profile and proximity accept it if the clause they cite for issuer data
  authentication has no validity check.

  Open point: the OpenID4VP and proximity profiles cite a different clause of
  ISO/IEC 18013-5 from the mdoc profile, with must, and do not say whether it
  covers time validity (C-101).  Is a relying party required to refuse an
  expired credential in all three profiles?

  Lookup: ISO/IEC 18013-5:2021 §9.1.2 and §9.3.1.

  Source:
  \href{https://standards.atlassian.net/wiki/spaces/OneID/pages/2947416082/Technical+Implementation+Guide}{One ID guide} §4.1.2, §4.2.
\end{theorem}
